% SOURCE_AUTHORITY: sga5_fr_workpass.tex, lines 14575--14603 (LF-normalized unit).
% SOURCE_SHA256: 791F4EFFC5E02832D5D77ED03518C8156D6F07E4C8238B03545DB93D883FBB28
Naturalmente, define-se da mesma forma uma classe de correspondência de Frobenius iterada $(\fr_X^\nu,(\Fr^*_{F/X})^\nu)$ para todo inteiro $\nu\geq0$; o isomorfismo
\[
(\Fr^*_{F/X})^\nu:(\fr_X^\nu)^*(F)\longrightarrow F
\]
é igual a
\[
(\Fr^*_{F/X})^{\nu-1}\circ (\fr_X^{\nu-1})^*(\Fr^*_{F/X})
\]
ou também a
\[
\Fr^*_{F/X}\circ\fr_X^*((\Fr^*_{F/X})^{\nu-1})
\]
para $\nu\geq1$.

Consideremos o caso particular em que $F$ é representável por um pré-esquema $V$ étale sobre $X$; então $(\fr_X)^*(F)=V^{(p/X)}$, $F(V)=\Hom_X(V,V)$,
\[
(\fr_X)^*(F)(V)=F(V^{(p/X)})=\Hom_X(V^{(p/X)},V),
\]
e $F(\Fr_{V/X})$ é a aplicação de $\Hom_X(V^{(p/X)},V)$ em $\Hom_X(V,V)$ que transforma um morfismo $\varphi$ em $\varphi\circ\Fr_{V/X}$; ora, $\Fr^*_{F/X}$ identifica-se com um morfismo de $V^{(p/X)}$ em $V$ que é a imagem de $\id_V$ pela aplicação $F(\Fr_{V/X})^{-1}$, inversa da precedente. Isto mostra que, para $F=V$ étale sobre $X$, tem-se
\[
\Fr^*_{V/X}=(\Fr_{V/X})^{-1}.
\]

Observe que, se o feixe $F$ é constante, é claro que
\[
(\fr_X)^*(F)=F,
\qquad
\Fr^*_{F/X}=\id_F.
\]
